\documentclass[12pt,a4paper]{article}
\usepackage[T1]{fontenc}
\usepackage{amsmath,amssymb,mathrsfs,tikz,times,pifont}
\usepackage{enumitem}
\usepackage{multicol}
\usepackage{lmodern}
\usetikzlibrary{trees}
\newcommand\circitem[1]{%
\tikz[baseline=(char.base)]{
\node[circle,draw=gray, fill=red!55,
minimum size=1.2em,inner sep=0] (char) {#1};}}
\newcommand\boxitem[1]{%
\tikz[baseline=(char.base)]{
\node[fill=cyan,
minimum size=1.2em,inner sep=0] (char) {#1};}}
\setlist[enumerate,1]{label=\protect\circitem{\arabic*}}
\setlist[enumerate,2]{label=\protect\boxitem{\alph*}}
%%%::::::by chnini ameur :::::::%%%
\everymath{\displaystyle}
\usepackage[left=1cm,right=1cm,top=1cm,bottom=1.7cm]{geometry}
\usepackage[colorlinks=true, linkcolor=blue, urlcolor=blue, citecolor=blue]{hyperref}
\usepackage{array,multirow}
\usepackage[most]{tcolorbox}
\usepackage{varwidth}
\usepackage{float} %pour utiliser l'option [H] qui force l'image à apparaître exactement à l'endroit où elle est placée dans le code.
\tcbuselibrary{skins,hooks}
\usetikzlibrary{patterns}
%%%::::::by chnini ameur :::::::%%%
\newtcolorbox{exa}[2][]{enhanced,breakable,before skip=2mm,after skip=5mm,
colback=yellow!20!white,colframe=black!20!blue,boxrule=0.5mm,
attach boxed title to top left ={xshift=0.6cm,yshift*=1mm-\tcboxedtitleheight},
fonttitle=\bfseries,
title={#2},#1,
% varwidth boxed title*=-3cm,
boxed title style={frame code={
\path[fill=tcbcolback!30!black]
([yshift=-1mm,xshift=-1mm]frame.north west)
arc[start angle=0,end angle=180,radius=1mm]
([yshift=-1mm,xshift=1mm]frame.north east)
arc[start angle=180,end angle=0,radius=1mm];
\path[left color=tcbcolback!60!black,right color = tcbcolback!60!black,
middle color = tcbcolback!80!black]
([xshift=-2mm]frame.north west) -- ([xshift=2mm]frame.north east)
[rounded corners=1mm]-- ([xshift=1mm,yshift=-1mm]frame.north east)
-- (frame.south east) -- (frame.south west)
-- ([xshift=-1mm,yshift=-1mm]frame.north west)
[sharp corners]-- cycle;
},interior engine=empty,
},interior style={top color=yellow!5}}
%%%%%%%%%%%%%%%%%%%%%%%
\usepackage{fancyhdr}
\usepackage{eso-pic}         % Pour ajouter des éléments en arrière-plan
% Commande pour ajouter du texte en arrière-plan
\usepackage{tkz-tab}
\AddToShipoutPicture{
    \AtTextCenter{%
        \makebox[0pt]{\rotatebox{80}{\textcolor[gray]{0.7}{\fontsize{5cm}{5cm}\selectfont PGB}}}
    }
}
\usepackage{lastpage}
\fancyhf{}
\pagestyle{fancy}
\renewcommand{\footrulewidth}{1pt}
\renewcommand{\headrulewidth}{0pt}
\renewcommand{\footruleskip}{10pt}
\fancyfoot[R]{
\color{blue}\ding{45}\ \textbf{2025}
}
\fancyfoot[L]{
\color{blue}\ding{45}\ \textbf{Prof:M. BA}
}
\cfoot{\bf
\thepage /
\pageref{LastPage}}
% Création du compteur pour les exercices
\newcounter{exercice}
\renewcommand{\theexercice}{\arabic{exercice}}  % Définit l'affichage du compteur en chiffres arabes

% Définir la commande \exo
\newcommand{\exo}{\refstepcounter{exercice}\textbf{Exercice \theexercice} }
\begin{document}
\renewcommand{\arraystretch}{1.5}
\renewcommand{\arrayrulewidth}{1.2pt}
\begin{tikzpicture}[overlay,remember picture]
    \node[draw=blue,line width=1.2pt,fill=purple,text=blue,inner sep=3mm,rounded corners,pattern=dots]at ([yshift=-2.5cm]current page.north) {\begingroup\setlength{\fboxsep}{0pt}\colorbox{white}{\begin{tabular}{|*1{>{\centering \arraybackslash}p{0.28\textwidth}} |*2{>{\centering \arraybackslash}p{0.2\textwidth}|} *1{>{\centering \arraybackslash}p{0.19\textwidth}|} }
                \hline
                \multicolumn{3}{|c|}{$\diamond$$\diamond$$\diamond$\ \textbf{Lycée de Dindéfélo}\ $\diamond$$\diamond$$\diamond$ } & \textbf{A.S. : 2025/2026}                                              \\ \hline
                \textbf{Matière: Mathématiques}                                                                                    & \textbf{Niveau : 1}\textbf{L} & \textbf{Date: 13/01/2026} & \textbf{} \\ \hline
                \multicolumn{4}{|c|}{\parbox[c]{10cm}{\begin{center}
                                                                  \textbf{{\Large\sffamily Td Polynomes}}
                                                              \end{center}}}                                                                                                        \\ \hline
            \end{tabular}}\endgroup};
\end{tikzpicture}

\vspace{3cm}
\fbox{\textbf{\exo}}

% --- IMAGE 6 : POLYNÔMES ---
\textbf{A) Compléter les phrases suivantes}
\begin{enumerate}
    \item Si $\alpha$ est racine d'un polynôme $P(x)$ alors $P(x)$ est factorisable par .............
    \item La forme factorisée d'un trinôme $ax^2 + bx + c$ dont les racines sont $x_1$ et $x_2$ est ...........
    \item Si $\alpha$ est une racine d'un polynôme de degré 4. Il existe $Q(x)$ de degré ....... tel que $P(x) = (x - \alpha)Q(x)$.
\end{enumerate}

\textbf{B) Répondre par vrai ou faux}
\begin{enumerate}
    \item $x^3 - y^3 = (x + y)(x^2 - xy + y^2)$
    \item Un polynôme de degré 3 admet 3 racines.
    \item Un polynôme qui a pour racine $\alpha$ et $\beta$ est factorisable par $(x - \alpha)(x - \beta)$.
\end{enumerate}
\fbox{\textbf{\exo}}
Soit le polynôme $P$ défini par : $P(x) = 2x^3 - 9x^2 - 8x + 15$.

\begin{enumerate}
    \item Calculer $P(5)$.
    \item Factoriser $P(x)$ par la méthode :
    \begin{enumerate}
        \item d'identification des coefficients ;
        \item de la division euclidienne ;
    \end{enumerate}
\end{enumerate}
\fbox{\textbf{\exo}}
Soit $P$ le polynôme défini par $P(x) = 3x^3 - x^2 - 8x - 4$.

\begin{enumerate}
    \item Vérifier que $2$ est une racine de $P$.
    \item Déterminer le polynôme $Q$ tel que $P(x) = (x - 2)Q(x)$.
    \item Factoriser $Q(x)$ puis en déduire une factorisation complète de $P(x)$.
    \item Résoudre dans $\mathbb{R}$ :
        a) $P(x) = 0$ \hfill b) $P(x + 2) = 0$ \hfill c) $P(x) < 0$
\end{enumerate}

\fbox{\textbf{\exo}}
Soit les polynômes $P(x) = -3x^3 - 2x^2 + 19x - 6$ et $T(x) = (1 - 3x)(2 - x)(-x - 3)$.

\begin{enumerate}
    \item 
    \begin{enumerate} 
    \item Calculer $P(1)$, $P(-3)$ et $P(0)$.
    
    \item Lequel des réels $1$, $-3$ et $0$ est une racine de $P$ ? Justifie
    \end{enumerate}
    \item Écrire $P(x)$ sous forme de produit de facteurs du premier degré. 
    
    \item 
    \begin{enumerate} 
    \item Donne en justifiant le degré du polynôme $T$. 
    \item Donner le monôme de plus haut degré de $T$. 
    
    \item Résoudre dans $\mathbb{R}$ l'équation $P(x) = 0$.
    
    \item Résoudre dans $\mathbb{R}$ l'inéquation $P(x) \leq 0$.
    \end{enumerate}
    \item Montrer que les polynômes $P$ et $T$ sont égaux.
\end{enumerate}

\fbox{\textbf{\exo}}

Soit le polynôme $f(x) = x^3 + ax^2 + bx + 6$ avec $a$ et $b$ sont deux nombres réels.

\begin{enumerate}
    \item Déterminer $a$ et $b$ sachant que $f(-2) = 0$ et $f(-1) = 8$.
    \item On pose $f(x) = x^3 - 2x^2 - 5x + 6$.
    \begin{enumerate}
        \item Calculer $f(1)$. Que peut-on en déduire ?
        \item Déterminer un polynôme $Q(x)$ tel que $f(x) = (x - 1)Q(x)$.
        \item Factoriser $Q(x)$ et en déduire une factorisation (complètement) de $f(x)$.
        \item Résoudre dans $\mathbb{R}$ :
        \begin{itemize}
            \item $f(x) = 0$
            \item $f(x) \geq 0$
        \end{itemize}
    \end{enumerate}
\end{enumerate}

\fbox{\textbf{\exo}}

\begin{enumerate}
    \item[I.] Définir les termes suivants : \\
    \textbf{Trinôme du second - Monôme de la variable $x$ - Polynôme}
    
    \item[II.] Résoudre dans $\mathbb{R}$ les équations et inéquations suivantes :
    \begin{enumerate}
        \item $x^2 - 6x = 0$ \quad  \hfill b. $-2x^2 + 3x + 2 = 0$
        \item $x^2 - 5x + 4 \leq 0$ \quad \hfill d. $-x^2 - x - 6 > 0$ 
    \end{enumerate}
\end{enumerate}

\fbox{\textbf{\exo}}

 On considère le polynôme $P$ défini par : $P(x) = x^3 - x^2 - 14x + 24$

\begin{enumerate}
    \item Calculer $P(1)$ ; $P(2)$ et $P(-3)$. 
    \item En déduire une racine de $P$. 
    \item Montrer qu'il existe un polynôme $Q$ tel que : $P(x) = (x - 2)Q(x)$ 
    \item Déterminer $Q(x)$. En déduire une factorisation complète de $P(x)$ 
    \item Résoudre dans $\mathbb{R}$ l'équation : $P(x) = 0$. En déduire les solutions dans $\mathbb{R}$ : \\
    $P(|x|) = 0$ 
    \vspace{0.3cm}
    \item Résoudre dans $\mathbb{R}$ l'équation : $P(x) \leq 0$ 
\end{enumerate}

\fbox{\textbf{\exo}}

\begin{enumerate}
    \item Trouver a et b tels que $2$ et $-1$ sont des racines de $P(x) = ax^3 + bx^2 - 5x - a$.
    
    \item En prenant les valeurs de a et b trouver dans 1), factoriser complètement $P(x)$.
    
    \item Résoudre dans $\mathbb{R}$ l'équation $P(x) = 0$. 
    
    \item En déduire les solutions des équations :
    \begin{enumerate}
        \item $P(x - 1) = 0$.
        \item $P(-|x|) = 0$. 
    \end{enumerate}
    
    \item Résoudre dans $\mathbb{R}$ l'inéquation $P(x) < 0$.
\end{enumerate}

\fbox{\textbf{\exo}} On considère le polynôme $P(x) = x^3 - 2x^2 - x + 2$

\begin{enumerate}
    \item Montrer que 1 est racine de $P(x)$. 
    
    \item En déduire une factorisation de $P(x)$.
    
    \item Résoudre l'équation $P(x) = 0$.
    
    \item Déduire les solutions des équations :
    \begin{enumerate}
        \item $(-x + 1)^3 - 2(-x + 1)^2 - (-x + 1) + 2 = 0$.
        \item $\dfrac{1}{(x+3)^3} - \dfrac{2}{(x+3)^2} - \dfrac{1}{x+3} + 2 = 0$.
    \end{enumerate}
    \item Résoudre l'inéquation $\dfrac{x^3 - 2x^2 - x + 2}{x^2 - x - 6} > 0$.
\end{enumerate}


\fbox{\textbf{\exo}}

On donne $P(x) = 2x^4 - 7x^3 - 2x^2 + 13x + 6$

\begin{enumerate}
    \item Calculer $P(0)$ et $P(-1)$ puis conclure. 
    \item Déterminer le polynôme $Q(x)$ tel que $P(x) = (x + 1)Q(x)$ 
    \item Dans la suite de l'exercice, on suppose que $Q(x) = 2x^3 - 9x^2 + 7x + 6$
    \begin{enumerate}
        \item Sachant que 2 est une racine de $Q(x)$, factoriser complètement $Q(x)$. 
        \item En déduire une factorisation de $P(x)$.
        \item Résoudre dans $\mathbb{R}$ l'équation $P(x) = 0$.
        \item Déduire de la question 3c) la résolution dans $\mathbb{R}$ de l'équation $P(x + 1) = 0$.
        \item Résoudre dans $\mathbb{R}$ l'inéquation $P(x) \leq 0$.
    \end{enumerate}
\item On pose : $f(x) = \dfrac{Q(x)}{-x^2 + 7x - 10}$
    \begin{enumerate}
        \item Déterminer le domaine de définition de $f$.
        \item Simplifier $f(x)$ puis résoudre dans $\mathbb{R}$ l'inéquation $f(x) > 0$.
    \end{enumerate}
\end{enumerate}

\fbox{\textbf{\exo}}

On considère le polynôme $P$ défini par : $P(x) = x^3 + 3x^2 - 4$.

\begin{enumerate}
    \item Montrer que 1 est une racine de $P$.
    \item Donner une factorisation de $P(x)$.
    \item Résoudre dans $\mathbb{R}$ :
    \begin{enumerate}
        \item l'équation $P(x) = 0$.
        \item l'inéquation $P(x) \geq 0$.
    \end{enumerate}
\end{enumerate}

\fbox{\textbf{\exo}}

On considère le polynôme $P$ défini par : $P(x) = -x^3 + 6x^2 - 9x + k$ où $k$ est un nombre réel.

\begin{enumerate}
    \item Déterminer la valeur du réel $k$ pour que 4 soit une racine de $P$.
    \item Pour la valeur de $k$ trouvée en 1) :
    \begin{enumerate}
        \item Factoriser $P$.
        \item Résoudre dans $\mathbb{R}$ l'équation $P(x) = 0$.
        \item Résoudre dans $\mathbb{R}$ l'inéquation $P(x) < 0$.
    \end{enumerate}
\end{enumerate}

\end{document}